% The abstract (PI 2026-11-24, LEDGER RW-9: the PI's draft, at most 180 words). Each clause is tested in
% experiments/FIGS/mlst_v2_tables.py (claims), which stops the build when a test is missing or fails. \changed marks
% what differs from the draft: the JetClass clause and the closing sentence, which the draft left open, and the scope
% of the mass clause (fine-tuned, the 17-class model with the mass regression is significantly less accurate).
% The draft has 171 words outside its two brackets, so the text below is at the limit of 180 without the closing
% sentence; checkpoint.tex prints that sentence and the clause with its two numbers under the abstract, uncounted.
Jet foundation models are often pretrained to classify simulated jets into hundreds of fine classes. Whether such
fine labels improve the pretrained model has not been tested in a controlled way, because published comparisons
change the training jets together with the labels. We pretrain the Sophon model on the same JetClass-II jets with
four nested label sets of \Kfull, \Kres, \Kmid{} and \Kcoarse{} classes, which progressively merge quark and lepton
flavors while keeping the decay topology. With coarser labels, a linear classifier on the frozen features identifies
the \Kfull{} classes less accurately, \VtwoAccCoarsePct\% against \VtwoAccFullPct\% for the coarsest and the finest
set, and decays that differ only in a merged flavor become much harder to separate. Random groupings into \Kcoarse{}
classes lose far less, so the loss depends on which classes are merged and not only on how many. After fine-tuning,
\changed{the coarsest set remains slightly worse on JetClass} and no significant
difference is seen on the top-tagging and quark--gluon benchmarks. An auxiliary jet-mass regression sharpens the mass
resolution of the frozen features at no significant cost in \changed{frozen-feature} classification.
